\chapter*{Introduzione}
\addcontentsline{toc}{chapter}{Introduzione}
\markboth{Introduzione}{Introduzione}

Le foreste sono una risorsa importante dal punto di vista ambientale ed economico, e gestirle in modo sostenibile richiede dati quantitativi sulla loro struttura, ad esempio quanti alberi e quante rocce ci sono in un'area boschiva, dove si trovano, la loro altezza e, soprattutto, qual è il diametro del fusto. Quest'ultimo parametro si misura per convenzione a circa 1,3~m dal suolo ed è indicato con la sigla DBH (\textit{Diameter at Breast Height}, diametro a petto d'uomo). Da queste grandezze si stimano il volume di legname, la biomassa e lo stato di salute del bosco, dati che nell'insieme prende il nome di inventario forestale~\cite{hamilton1975}.

Tradizionalmente l'inventario viene realizzato a mano. Alcuni operatori percorrono l'area interessata e misurano gli alberi uno per uno con strumenti dendrometrici. Il metodo è affidabile, ma lento e costoso in termini di personale, difficile da applicare su superfici estese e da ripetere ogni volta che si vuole seguire l'evoluzione del bosco. Questo approccio può risultare carente anche in termini di sicurezza del personale, dato che alcuni ambienti boschivi possono rivelarsi particolarmente pericolosi per la loro stessa conformazione geologica. Per questo negli ultimi decenni si è cercato di automatizzare la raccolta dei dati.

Un contributo decisivo è arrivato dal LiDAR (\textit{Light Detection and Ranging}), un sensore che emette impulsi laser e ne misura il tempo di ritorno, ricavando la distanza dei punti colpiti. Il risultato è una nuvola di punti (\textit{Point Cloud}) che descrive la geometria dell'ambiente e che si ottiene in tempi brevi, motivo per cui il LiDAR è diventato uno strumento molto usato nello studio della struttura forestale~\cite{lim2003}. Le prime applicazioni erano da aereo o da satellite, questo permette indubbiamente di coprire aree vaste, ma osservano il bosco dall'alto, e la chioma degli alberi intercetta una parte consistente degli impulsi. I fusti restano in gran parte nascosti e il loro diametro, il dato più utile per l'inventario, è difficile da stimare con precisione. Per osservare i tronchi il sensore va dunque inevitabilmente portato a terra, dentro il bosco.

Oggi una soluzione diffusa è far attraversare il bosco a un operatore con uno scanner LiDAR portatile, per esempio montato su uno zaino~\cite{gennai2026}. Rispetto alle misure manuali i tempi si riducono molto e si ottiene una descrizione tridimensionale dettagliata, ma serve ancora una persona fisica che percorra tutta l'area. Il passo successivo è far trasportare il sensore da un robot~\cite{mattamala2025}. Il terreno forestale però è irregolare, spesso in pendenza e pieno di radici, rocce e rami caduti, e per veicoli a ruote o a cingoli è difficile da affrontare. Per questo i robot quadrupedi rappresentano una scelta particolarmente adatta. Un esempio è il robot Spot di Boston Dynamics, impiegato a questo scopo nella tesi magistrale~\cite{gennai2026} da cui parte questo lavoro.

Avere un robot che porta il sensore non basta. Perché l'acquisizione delle informazioni sia davvero automatica, il robot deve decidere da solo dove muoversi per mappare l'area assegnata, senza conoscerla in anticipo e senza un operatore che lo guidi. Serve quindi un algoritmo di esplorazione autonoma, che a ogni passo scelga dove proseguire in base a ciò che il sensore ha già osservato. Lo studio e il confronto di algoritmi di questo tipo sono l'oggetto di questa tesi.

Esplorare significa costruire la mappa di un ambiente sconosciuto usando solo le misure raccolte durante il movimento. A ogni passo il robot deve risolvere due problemi diversi: dove andare, cioè quale destinazione scegliere, e come arrivarci, cioè quale percorso seguire nella parte di mappa già nota. Il secondo quesito è un classico problema di pianificazione del percorso mentre il primo rappresenta esattamente il compito che dovrà essere svolto dalla strategia di esplorazione ed è quello su cui si concentra questo lavoro. Un approccio molto diffuso si basa sulle frontiere, definite da Yamauchi~\cite{yamauchi1997} come il confine tra lo spazio libero già osservato e quello ancora sconosciuto. Un'alternativa viene dagli algoritmi di copertura come la decomposizione \textit{boustrophedon}~\cite{choset2000}, che percorrono l'area a fasce parallele.

Il lavoro ha richiesto per prima cosa la realizzazione di un simulatore bidimensionale di un bosco, con un sensore LiDAR simulato e una mappa a griglia di occupazione~\cite{elfes1989} aggiornata a ogni scansione. Il modello in due dimensioni ha una giustificazione precisa: a quota fissa il LiDAR osserva una sezione orizzontale del bosco, in cui i tronchi appaiono come cerchi e le rocce come poligoni irregolari. La simulazione permette inoltre di generare molti boschi diversi in modo riproducibile e di confrontare le strategie a parità di ambiente, cosa difficile con prove sul campo. Su questa base sono state implementate tre strategie differenti: una casuale, usata come riferimento, una basata sulle frontiere e una a serpentina ispirata agli algoritmi di copertura. Tutte usano lo stesso sensore, la stessa mappa e lo stesso algoritmo di pianificazione, quello di Dijkstra~\cite{dijkstra1959}, e differiscono solo nella scelta della destinazione, così che le differenze osservate dipendano unicamente dalla strategia. Il confronto è stato condotto su 50 boschi generati casualmente, con due diversi punti di partenza, misurando la copertura raggiunta e la distanza percorsa. I casi in cui l'esplorazione si interrompe sono stati analizzati nel dettaglio, e da questa analisi è nata una variante della strategia basata sulle frontiere, confrontata con la versione originale sugli stessi boschi dopo averne tarato il parametro su un insieme di boschi distinto.

Rispetto alla tesi magistrale~\cite{gennai2026}, che ha sviluppato algoritmi di esplorazione per Spot, il problema è affrontato qui in un ambiente simulato e semplificato, per concentrarsi sul confronto tra le strategie. Su indicazione del relatore è stato riutilizzato il modello del sensore LiDAR sviluppato in quel lavoro. Tutte le altre parti del sistema sono state sviluppate in questa tesi.

La strategia basata sulle frontiere è risultata la più efficiente, ma anche la più soggetta a blocchi, situazioni in cui il robot resta circondato da celle non percorribili e deve fermarsi. La strategia casuale non si blocca mai con partenza dal centro e raggiunge la copertura più alta, ma percorre in media circa 360~m, contro i circa 150~m della strategia basata sulle frontiere e i 190~m della serpentina, che si colloca in posizione intermedia anche per numero di blocchi. Partendo da un angolo dell'area le conclusioni non cambiano. I blocchi hanno una causa comune, legata al modo in cui il robot verifica il percorso e agli ostacoli che il sensore vede di striscio e scopre solo da vicino. La variante proposta penalizza le destinazioni che si trovano in passaggi stretti, tra due ostacoli o tra un ostacolo e il bordo dell'area. Tramite questa modifica nessuna esplorazione si interrompe sotto il 90\% di copertura, contro le 10 della versione originale, e la distanza percorsa aumenta di circa il 4\%.

Il Capitolo~\ref{cap:metodologia} descrive il modello dell'ambiente e del sensore, la costruzione della mappa, la pianificazione del percorso e le strategie di esplorazione, compresa la variante proposta. Il Capitolo~\ref{cap:implementazione} presenta l'implementazione in Python e gli strumenti per eseguire i test e visualizzare l'esplorazione. Il Capitolo~\ref{cap:risultati} riporta i risultati del confronto e l'analisi dei blocchi. Le Conclusioni riassumono il lavoro e indicano alcuni possibili sviluppi.
